\documentclass[11pt,a4paper]{article}
\usepackage[utf8]{inputenc}
\usepackage[T1]{fontenc}
\usepackage{libertine}
\usepackage[margin=0.9in]{geometry}
\usepackage{hyperref}
\hypersetup{colorlinks=true,urlcolor=blue}
\setlength{\parindent}{0pt}
\setlength{\parskip}{0.85em}
\pagestyle{empty}

\begin{document}

\begin{flushright}
Trento, 28 ottobre 2026
\end{flushright}

\textbf{Oggetto: Candidatura per il tirocinio IT Application Analyst -- Artificial Intelligence}

Gentile Ufficio Selezione FERCAM,

sono Filippo Lollato, studente magistrale in Artificial Intelligence Systems all'Università di Trento e laureato triennale in Informatica. Desidero candidarmi per il tirocinio IT Application Analyst -- Artificial Intelligence presso la sede di Bolzano. Mi interessa in particolare il taglio pratico del ruolo: capire dove l'AI può essere utile nei processi aziendali e contribuire a integrare strumenti che funzionino in modo affidabile nel contesto esistente.

Nel progetto universitario ASA Autobots ho co-guidato un team di due persone nello sviluppo di un workflow per le consegne che integra agenti BDI, un coordinatore basato su LLM e un planner PDDL. Ho contribuito all'architettura e all'integrazione dei componenti e delle interfacce tra servizi. Questa esperienza mi ha dato modo di affrontare aspetti concreti dell'impiego di agenti e modelli linguistici all'interno di un sistema più ampio.

La mia tesi triennale è diventata \textit{incin}, un progetto personale in Rust che esplora come controlli strutturali e feedback del compilatore possano rendere il codice di machine learning più affidabile. In un progetto universitario di NLP ho inoltre realizzato una pipeline di addestramento e valutazione con Python, PyTorch e Hugging Face. Nel mio ruolo di tutor universitario in Rust, Programmazione Avanzata e Compilatori, ho imparato a spiegare concetti tecnici con chiarezza e a supportare le persone nella risoluzione dei problemi.

Vorrei mettere queste esperienze al servizio delle attività descritte nell'annuncio, dalla valutazione di strumenti AI alla documentazione e all'integrazione con le soluzioni aziendali, continuando a imparare a stretto contatto con il team IT e con le funzioni aziendali. Sono interessato al percorso di tirocinio di 4--6 mesi indicato e disponibile ad approfondirne modalità e calendario.

Allego il mio curriculum e resto volentieri a disposizione per un colloquio.

Cordiali saluti,\\
Filippo Lollato\\
\href{mailto:contact@lollato.dev}{contact@lollato.dev} \enskip\textbullet\enskip \href{https://lollato.dev}{lollato.dev}

\end{document}
